\chapter{Introduction}
\label{chap:intro}

Communication is a fundamental aspect of life on earth, essential for conveying
information, coordinating behaviors, and establishing social relationships.
From intricate visual displays to complex vocalizations, communication plays a
crucial role in the survival and reproduction of species across the tree of
life. In humans, communication takes on a myriad of forms, from spoken language
over gestures and facial expressions to written text. The study of animal
communication offers insights into the ecology and evolution of
species, providing a window into the underlying mechanisms and functions of
different communication signals used in various contexts. At the same time, it
provides a platform that could help us understand how communication generally
and our language specifically evolved and how it is produced and perceived in
the brain.

\ornament

Electric fish offer a unique opportunity to study animal communication and
social behavior but are rarely used as a model system outside the field of
Neuroethology. Beyond utilizing their electric field for active sensing,
wave-type electric fish have been used extensively to study the physiological
underpinnings of the production and perception of communication signals. The
most frequently observed, and hence studied, electrocommunicaion signal are
so-called chirps. During a chirp, electric fish briefly and transiently
increases the frequency with which it discharges its electric organ, a
phenomenon found in many species
\parencite{petzoldCoadaptationElectricOrgan2016,
smithEvolutionElectricCommunication2016,
freilerElectrocommunicationSignalsAggressive2022,
zhouStructureSexualDimorphism2006, hoSexDifferencesElectrocommunication2010}.

Understanding chirps necessitates a introduction of the \acrfull{EOD} in
wave-type weakly electric fish. The EOD, a quasi-sinusoidal waveform similar to
a constant-pitch tone in audio, features a fundamental frequency
\acrshort{EOD$f$} and harmonics. Its amplitude decays dependent on the fish's
orientation, body curvature, and distance relative to the recording electrode
\parencite{bendaPhysicsElectrosensoryWorlds2020}. While EOD$f$ differs among
individuals, it's generally stable, with variations mainly due to water
temperature changes \parencite{moortgatSubmicrosecondPacemakerPrecision1998}.
Frequencies range from less than \SI{50}{\hertz} to over \SI{2}{\kilo\hertz},
with each species occupying a characteristic frequency range
\parencite{bendaPhysicsElectrosensoryWorlds2020,
evansTaxonomicRevisionDeep2017a}. Despite its simplicity, the EOD waveform can
reveal species-specific traits under standardized recording conditions.

Wave-type weakly electric fish can modulate the frequency of their EOD on time
scales of milliseconds to minutes to produce various communication signals
\parencite{zakonEODModulationsBrown2002a}. In most cases, the chirps of
\textit{A.~leptorhynchus} are \enquote{simple} up-modulations of the EOD$f$
that last between \SI{20}{\milli\second} and \SI{500}{\milli\second} and can
reach between \SI{20}{\hertz} and \SI{500}{\hertz} above the baseline EOD$f$.
From now on, when I refer to a chirp as \enquote{short} or \enquote{long}, I
will refer to the duration of the chirp. When mentioning the properties
\enquote{small} or \enquote{large}, I will refer to the frequency up-modulation
from baseline. Chirps can come in many shapes: Small and short chirps are
usually approximately shaped like a Gaussian as illustrated in figure
\ref{fig:multifish_chirps} A. If they become larger, this Gaussian can be
followed by a brief undershoot. In contrast, the shape of long and large chirps
usually vaguely resembles a square wave as in figure \ref{fig:multifish_chirps}
E. If we now take a look at the signal that underlies this change in frequency,
so the signal that we could record, we can see, that the third feature of a
chirp, in addition to the duration and the frequency change, is the change in
amplitude. For small chirps (figure \ref{fig:multifish_chirps} B) this trough
in amplitude is usually not as pronounced as it is for larger chirps (figure
\ref{fig:multifish_chirps} F).

\begin{figure}
  \centering
  \includegraphics[width=\textwidth]{../figures/multifish_chirps.pdf}
  \caption{
    Overview of how chirps as changes in the EOD$f$ of a fish can be recorded
    in different settings if other fish are around, simulated from the
    instantaneous frequency evolutions from chirps of real fish. Chirps are
    defined as a brief elevation of the EOD$f$ of one fish. They can be
    relatively small and short (A) or long with a high up-modulation (E). If
    only the chirping fish would be recorded in a chirp-chamber-like setup, a
    small chirp is already hardly visible by just the EOD alone (B) while a
    larger chirp might appear a bit clearer (F). In the field however, the EODs
    of multiple fish superimpose, which makes chirp detection but also chirp
    assignment much harder without a frequency spectrum (C, G). Because of
    this, we need a \acrfull{STFT} to convert the recordings to spectrograms,
    which allows an approximation of the underlying discharge frequencies (D,
    H). As this is just an approximation, the resolvability of chirps and the
    separability of fish forms a trade-off: The temporal resolution of the
    spectrogram, if sufficient to separate fish, is not sufficient to resolve a
    full chirp (D) but rather produces an \enquote{artifact} caused by the
    presence of the chirp which is usually much longer than the underlying
    chirp, particularly for small chirps. For long chirps (H), this is not so
    much of a problem.
  }
  \label{fig:multifish_chirps}
  \vspace{0.5cm}
\end{figure}

Chirps offer a unique opportunity to link advances in animal behavior to the
underlying neural mechanisms: They are extensively studied in controlled
conditions and the endocrinological as well as neurological mechanisms to
produce and perceive them are comparatively well understood
\parencite{hupeEffectDifferenceFrequency2008,
zupancElectrogenesisElectroreceptionOverview2005,
smithSerotonergicActivation5HT1A2008, dunlapGlucocorticoidReceptorBlockade2011,
dulkaTestosteroneModulatesFemale1994, kellerControlPacemakerModulations1991,
kraheNeuralMapsElectrosensory2014, marsatCellularCircuitProperties2012,
dunlapVocalElectricFish2021, bendaSpikeFrequencyAdaptationSeparates2005,
bendaSynchronizationDesynchronizationCodeNatural2006}. But the role of chirps
produced in various naturalistic behavioral contexts is still mysterious.
Functionally, chirps as a communication signal have been attributed to roles in
courtship as well as signaling dominance and submission, the latter two seeming
contradictory \parencite{batistaNonsexbiasedDominanceSexually2012,
hagedornCourtSparkElectric1985a, englerDifferentialProductionChirping2001,
henningerStatisticsNaturalCommunication2018}. On the other hand, recent work
suggested that chirps might not serve as a means of conventional communication
at all but as a form of autocommunication to increase the capabilities of the
active electric sense in social contexts \parencite{obotiWhyBrownGhost2023}.
Except for the work by \textcite{henningerStatisticsNaturalCommunication2018},
most work underlying proposals for various functional roles of chirps rely
heavily on data recorded in artificial conditions and none of them provide
causal evidence. But to correctly interpret physiological results holistically,
we must include a species ecology and evolution
\parencite{tinbergenAimsMethodsEthology1963}. A dramatic example is the work by
\textcite{henningerStatisticsNaturalCommunication2018}, who showed that the
stimulus regimes commonly used in controlled conditions to study the neural
mechanisms underlying electric communication are not representative of the
natural environment. I argue that before probing for potential explanations for
a specific behavior in highly controlled conditions, we should aim to first
\emph{observe} correlations that suggest causal relationships to generate
\emph{informed} hypotheses before we test them with controlled experiments. In
other words: We need to observe which fish chirps at what time and during which
kind of social interactions in a naturalistic context. 

The biggest challenge preventing these approaches to studying chirping is that
chirps are hard to detect automatically in behaviorally relevant settings.
Currently, chirps can be automatically detected in cases where the fish are
separated from each other so that their electric fields do not superimpose
\parencite{lehotzkySupervisedLearningAlgorithm2023,
eskeEffectUrethaneMS2222023, fieldJARChirpsGymnotiform2019}. But studies where
fish are kept together are much more informative about their natural behavior
and often uncover behaviors that challenge common assumptions from previous
literature \parencite{henningerStatisticsNaturalCommunication2018,
fortuneSpookyInteractionDistance2020a}. In these cases, the electric fields of
the fish superimpose and the chirps are much harder to detect
\parencite{obotiWhyBrownGhost2023,henningerStatisticsNaturalCommunication2018}.
This is either not studied at all, or if so involves a lot of manual
annotation, which is a labor-intensive and error-prone process and not feasible
for large datasets, such as the ones recorded in the field. But it is not the
problem that we do not have the tools to detect chirps in these cases. As an
example, smartphone apps can automatically recognize bird songs and classify
the species just from comparatively poor-quality phone recordings
\parencite{kahlBirdNETDeepLearning2021}. Despite this, many studies still rely
on manually labeling communication signals of all kinds of species. Also,
software such as
\href{http://www.mackenziemathislab.org/deeplabcut}{\texttt{DeepLabCut}} and
\href{https://sleap.ai/}{\texttt{SLEAP}} enables automated tracking of
individual body parts of most animals \parencite{pereiraSLEAPDeepLearning2022a,
lauerMultianimalPoseEstimation2022}. This provides the possibility of high
throughput in behavioral experiments and frees time that can be allocated to,
for example, recording more data. Algorithms to automate the conversion of
unstructured data (such as video, audio or electric recordings) into structured
data (such as tables to do statistics on) are widely available, they just need
to be embedded into software packages to make them accessible to researchers.

In this thesis, I want to provide a demonstration of such a fully automated
pipeline. Here, I aim to implement a first proof of concept for fully automated
chirp detection in large datasets recorded in the field for \textit{A.
leptorhynchus}, although the pipeline can be extended to any chirping species
of wave-type electric fish. As you are about to see, this does not require
inventing new technologies but simply connecting existing, openly available and
well documented tools into a data pipeline to fully automate the detection of
chirps in large datasets for quantitative research in electrocommunication.
This will not only enable to observe and hence analyze chirps in field
recordings but could also be an important building block for electric fish
monitoring.

To understand the features that are useful to detect chirps, I will first
introduce the physical properties of the \acrfull{EOD} of weakly electric and
how chirps change it in the \hyperref[chap:background]{next} chapter. We will
then transition to a brief introduction on neural networks and how they can
applied to the problem of chirp detection. In the
\hyperref[chap:methods]{following} chapter, I will introduce the training data
generation pipelines for machine learning models, how the models are trained
and how their performance is evaluated. I will also explain how the detector
can be applied to a large dataset of electric recordings and how I analyzed
behaviors associated with chirps. In the \hyperref[chap:results]{results}, we
will look at the full architecture of the chirp detection pipeline, the
performance of the neural networks and the results of the analysis of the large
dataset of weakly electric fish recordings. And finally, in the
\hyperref[chap:discussion]{discussion} chapter, we will discuss the
implications of the results and the limitations of the pipeline. We will also
discuss how the pipeline can be improved and how it can be applied to other
species and other tasks beyond analyzing animal communication, such as
environmental monitoring. For now, let's start with the fundamentals: How do we
record the electric fields of weakly electric fish?

